\begin{theorem}[Хватал, 1977]
Для любого дерева \(T_n\) на \(n\) вершинах и \(m\ge1\)
\[
R(T_n,K_m)=(m-1)(n-1)+1 .
\]
Здесь ищется дерево \(T_n\) цвета 1 или клика \(K_m\) цвета 2. Ответ не зависит от формы дерева!
\end{theorem}
\textit{Нижняя оценка: \(R(T_n,K_m)\ge(m-1)(n-1)+1\).}
Раскрасим рёбра \(K_{(m-1)(n-1)}\): разобьём вершины на \(m-1\) групп по \(n-1\) вершине,
рёбра внутри групп — цвета 1, между группами — цвета 2. Компоненты связности
графа цвета 1 имеют по \(n-1\) вершине, поэтому связного \(T_n\) цвета 1 нет. Граф
цвета 2 — полный \((m-1)\)-дольный: среди любых \(m\) вершин две в одной доле, они не
соединены ребром цвета 2, так что \(K_m\) цвета 2 нет.
\begin{graphillustration}{Экстремальная раскраска при $m=3$, $n=4$: две группы по 3 вершины. Внутри групп цвет 1 (сплошные) — компоненты из $n-1=3$ вершин, $T_4$ не помещается; между группами цвет 2 (штриховые) — двудольный граф без треугольника.}{fig:12-chvatal}
\begin{tikzpicture}[font=\small]
\node[gi vertex] (a1) at (0,1.4) {};
\node[gi vertex] (a2) at (-0.8,0) {};
\node[gi vertex] (a3) at (0.8,0) {};
\node[gi vertex] (b1) at (4,1.4) {};
\node[gi vertex] (b2) at (3.2,0) {};
\node[gi vertex] (b3) at (4.8,0) {};
\foreach \x in {a1,a2,a3} \foreach \y in {b1,b2,b3} \draw[gi c2] (\x) -- (\y);
\draw[gi c1] (a1) -- (a2) -- (a3) -- (a1);
\draw[gi c1] (b1) -- (b2) -- (b3) -- (b1);
\node[gi group,fit=(a1)(a2)(a3),label=below:{$K_{n-1}$}] {};
\node[gi group,fit=(b1)(b2)(b3),label=below:{$K_{n-1}$}] {};
\node[gi note,anchor=west] at (5.6,0.7) {$(m-1)(n-1)=6$ вершин:\\нет $T_4$ цвета 1,\\нет $K_3$ цвета 2};
\end{tikzpicture}
\end{graphillustration}


\begin{lemma}\label{lem:trees}
Пусть \(m>1\), \(H\) — граф с \(v(H)\ge(m-1)(n-1)+1\) и \(\alpha(H)\le m-1\). Тогда \(H\)
содержит (как подграф) любое дерево на \(n\) вершинах.
\end{lemma}
\begin{proof}
Индукция по \(n\) при фиксированном \(m\). База \(n=1\): в \(H\) есть хотя бы одна вершина.

Переход \(n-1\to n\). Пусть \(u\) — лист \(T_n\), \(a\) — его сосед, \(T_{n-1}=T_n-u\).
Возьмём максимальное независимое множество \(U\) графа \(H\): \(|U|=\alpha(H)\le m-1\). Тогда
\[
v(H-U)\ge(m-1)(n-1)+1-(m-1)=(m-1)(n-2)+1,\qquad \alpha(H-U)\le\alpha(H)\le m-1,
\]
и по предположению индукции \(H-U\) содержит копию \(T_{n-1}\). Пусть \(a'\) — образ \(a\).
Множество \(U\cup\{a'\}\) не независимо (иначе \(U\) не максимально), поэтому \(a'\) смежна
с некоторой \(x\in U\). Вершина \(x\) не входит в копию \(T_{n-1}\subseteq H-U\);
присоединив её к \(a'\), получаем копию \(T_n\).
\end{proof}
\begin{graphillustration}{Шаг индукции в лемме: копия $T_{n-1}$ найдена в $H-U$; вершина $a'$ (образ соседа удалённого листа) смежна с некоторой $x\in U$ по максимальности $U$. Ребро $a'x$ (толстое) достраивает $T_n$.}{fig:12-step}
\begin{tikzpicture}[font=\small]
\node[draw,rounded corners,minimum width=5.2cm,minimum height=2.4cm] (HU) at (0,0) {};
\node[gi note,anchor=north west] at (-2.55,1.15) {$H-U$};
\node[gi vertex] (t1) at (-1.6,-0.2) {};
\node[gi vertex] (t2) at (-0.6,0.4) {};
\node[gi vertex] (t3) at (-0.6,-0.7) {};
\node[gi vertex] (t4) at (0.5,0.4) {};
\node[gi selected] (a) at (1.6,0.4) {$a'$};
\draw[gi edge] (t1) -- (t2) -- (t4) -- (a); \draw[gi edge] (t1) -- (t3);
\node[draw,rounded corners,minimum width=1.4cm,minimum height=2.4cm] (U) at (4.2,0) {};
\node[gi note,anchor=north] at (4.2,1.6) {$U$, $|U|\le m-1$};
\node[gi vertex] (x) at (4.2,0.6) {$x$};
\node[gi vertex] (y) at (4.2,-0.6) {};
\draw[gi match] (a) -- (x);
\end{tikzpicture}
\end{graphillustration}


\textit{Верхняя оценка.} Рассмотрим раскраску \(K_N\), \(N=(m-1)(n-1)+1\), и граф \(H\)
рёбер цвета 1. Если нет \(K_m\) цвета 2, то в \(H\) нет независимого множества размера
\(m\): \(\alpha(H)\le m-1\) (при \(m=1\) утверждение тривиально: \(N=1\) и \(K_1\) есть). По
лемме~\ref{lem:trees} \(H\) содержит \(T_n\), т.\,е. есть \(T_n\) цвета 1. \qed

\begin{example}
\(R(P_3,K_3)=2\cdot2+1=5\), \(R(K_{1,3},K_3)=R(P_4,K_3)=7\): звезда и путь на 4 вершинах
дают одно и то же число.
\end{example}
